% ============================================================================
% Taller 1 — Componentes, descomposicion y estacionariedad (Modulo I)
% Series de Tiempo (20948) · Corte I · Universidad El Bosque · 2026-II
%
% PLANTILLA DE ENTREGA. Se entrega el PDF compilado, por Brightspace.
%
% Como usarla:
%   1. Rellena la ficha de la portada. El digito de verificacion es lo primero
%      que hay que comprobar: un ADF perfectamente calculado sobre la serie
%      equivocada se ve identico a uno correcto.
%   2. Rellena la hoja de cifras de la pagina 2 ANTES de escribir prosa.
%   3. Responde cada literal donde dice \marcador{...}. Los marcadores en gris
%      que queden sin sustituir se ven en el PDF: esa es su unica funcion.
%   4. Compila con pdflatex (o subela a Overleaf tal cual: no usa ningun
%      paquete raro ni ningun archivo externo) y entrega el PDF.
%
% Nombre del archivo:  T1_Apellido_XXX.pdf   (XXX = tus tres ultimos digitos)
% ============================================================================
\documentclass[11pt,a4paper,oneside]{article}

\usepackage[utf8]{inputenc}
\usepackage[T1]{fontenc}
\usepackage[spanish,es-nodecimaldot]{babel}
\usepackage{lmodern}
\usepackage{microtype}
\usepackage[a4paper,top=2.3cm,bottom=2.3cm,left=2.4cm,right=2.4cm]{geometry}

\usepackage{graphicx}
\usepackage{booktabs}
\usepackage{array}
\usepackage{tabularx}
\usepackage{caption}
\usepackage{enumitem}
\usepackage{fancyhdr}
\setlength{\headheight}{14pt}
\usepackage{titlesec}
\usepackage{amsmath,amssymb}
\usepackage{xcolor}
\usepackage{listings}
\usepackage{float}
\usepackage{hyperref}

% --- Colores institucionales (los mismos del material de Series de Tiempo) ---
\definecolor{ubosque}{RGB}{1,40,32}
\definecolor{naranja}{RGB}{255,102,0}
\definecolor{grisclaro}{RGB}{240,240,240}
\definecolor{grismarcador}{RGB}{130,130,130}

\hypersetup{colorlinks=true, linkcolor=ubosque, urlcolor=blue!70!black,
            pdftitle={Taller 1 - Series de Tiempo 20948}}

\pagestyle{fancy}
\fancyhf{}
\fancyhead[L]{\small\textit{Taller 1 --- Series de Tiempo (20948)}}
\fancyhead[R]{\small\textit{\leftmark}}
\fancyfoot[C]{\thepage}
\fancyfoot[R]{\small Universidad El Bosque}
\renewcommand{\headrulewidth}{0.4pt}
\renewcommand{\footrulewidth}{0.3pt}

\setlength{\parskip}{0.5\baselineskip}
\setlength{\parindent}{0pt}
\emergencystretch=3em
\titleformat{\section}{\normalfont\large\bfseries\color{ubosque}}{\thesection.}{0.5em}{}
\titleformat{\subsection}{\normalfont\normalsize\bfseries}{}{0em}{}
\captionsetup{font=small,labelfont=bf,labelsep=period}

\lstset{basicstyle=\ttfamily\footnotesize, breaklines=true, frame=single,
        rulecolor=\color{grisclaro!60!black}, columns=flexible,
        showstringspaces=false, xleftmargin=0pt, framexleftmargin=3pt}

% --- Los tres comandos de la plantilla ---
% Un hueco por rellenar. Se ve en gris y entre corchetes: si queda uno sin
% sustituir, salta a la vista en el PDF entregado.
% Ojo con la forma del color: `\textcolor{}{}` (comando) y NO
% `{\color{} ...}` (declaracion). Dentro de una celda de parrafo —las
% columnas X de tabularx— la declaracion arranca el parrafo con un whatsit y
% hunde su primera linea: la etiqueta de la fila queda arriba y su casilla
% media linea mas abajo. Aislado en un caso minimo: `{\itshape[x]}` alinea,
% `{\color{gris}[x]}` no. (Lección heredada de la plantilla de Estadistica
% Espacial; no cambiar sin reproducir el caso minimo.)
\newcommand{\marcador}[1]{\textcolor{grismarcador}{\textit{[#1]}}}

\newcommand{\tarea}[4]{%
  \section[#1 · #2]{#1 · #2 \normalfont\small\color{grismarcador}(#3 \% del escrito · #4 palabras)}%
  \markboth{#1 · #2}{}}

% Un literal: la letra en negrita y el rotulo corto. El enunciado COMPLETO
% esta en el taller; aqui solo va el rotulo, para no gastar dos paginas en
% copiar lo que el estudiante ya tiene delante.
\newcommand{\literal}[2]{\vspace{0.35\baselineskip}\noindent\textbf{(#1)}~\textit{#2}\par}

\begin{document}

% ============================================================================
% PORTADA — la ficha de variante
% ============================================================================
\thispagestyle{empty}
\begin{center}

{\large\textbf{UNIVERSIDAD EL BOSQUE}}\\[0.15cm]
{\normalsize Facultad de Ciencias --- Matemáticas y Ciencia de Datos}\\[0.1cm]
{\normalsize Series de Tiempo (20948) --- 2026-II --- Corte I}\\[0.9cm]

\rule{\textwidth}{1.2pt}\\[0.25cm]
{\LARGE\textbf{Taller 1}}\\[0.15cm]
{\large\textit{Componentes, descomposición\\ y estacionariedad}}\\[0.25cm]
\rule{\textwidth}{1.2pt}\\[0.9cm]

{\large\textbf{Trabajo individual}}\\[0.15cm]
{\normalsize Escrito 60\,\% \quad · \quad Defensa oral 40\,\%}\\[1.0cm]

\end{center}

\begin{tabular}{rl}
\textbf{Nombre completo:}          & \marcador{tu nombre} \\[0.25cm]
\textbf{Número de documento:}      & \marcador{completo} \\[0.25cm]
\textbf{Tus tres últimos dígitos:} & \marcador{000 a 999} \quad $\rightarrow$ \quad
                                     \textbf{variante} \marcador{la misma cifra} \\[0.25cm]
\textbf{Tu serie no estacionaria:} & número \marcador{NN} \quad con $n =$ \marcador{n} observaciones \\[0.25cm]
\textbf{Tu serie estacional:}      & número \marcador{NN} \quad con $n =$ \marcador{n} observaciones \\[0.25cm]
\textbf{Dígito de verificación:}   & \marcador{la suma que da la tira del taller} \\[0.25cm]
\textbf{Fecha de entrega:}         & \textbf{martes 25 de agosto de 2026, 14:00} \\
\end{tabular}

\vspace{0.9cm}

\noindent\fcolorbox{ubosque}{grisclaro}{%
\begin{minipage}{0.95\textwidth}
\small
\textbf{Antes de escribir nada, comprueba el dígito de verificación.} Es la suma de los valores de
tu serie, redondeada. Si la tuya no da exactamente esa cifra, no estás trabajando con tu variante,
y un ADF perfectamente calculado sobre la serie equivocada se ve idéntico a uno correcto. Todo lo
que construyas encima estaría mal sin que nada avise.
\end{minipage}}

\vspace{0.5cm}

\noindent\fcolorbox{grisclaro!60!black}{white}{%
\begin{minipage}{0.95\textwidth}
\small
\textbf{Cómo se entrega.} Un solo PDF por \textbf{Brightspace}, compilado con \LaTeX{} a partir de
esta plantilla. Nombre del archivo: \texttt{T1\_Apellido\_XXX.pdf}, donde \texttt{XXX} son tus tres
últimos dígitos. Máximo \textbf{11 páginas} sin contar el anexo. Las tablas y figuras no cuentan
para los límites de palabras de cada tarea.

\smallskip
\textbf{El enunciado no está aquí.} Vive en el taller publicado, con tu variante resuelta al
escribir tu documento. Esta plantilla solo recoge las respuestas.
\end{minipage}}

\vfill
{\footnotesize\color{grismarcador}\today}

\newpage

% ============================================================================
% HOJA DE CIFRAS
% ============================================================================
\section{Hoja de cifras}
\markboth{Hoja de cifras}{}

Todas las cifras que el taller pide, juntas. \textbf{Rellénala antes de escribir prosa}: es lo que
permite que la discusión de cada tarea hable de números que ya están decididos, y es lo primero que
se revisa al calificar. Usa punto decimal y cuatro decimales donde la cifra los tenga.

\vspace{0.3cm}
\renewcommand{\arraystretch}{1.25}
\begingroup\setlength{\parskip}{0pt}
\begin{tabularx}{\textwidth}{@{}l X r@{}}
\toprule
\textbf{Tarea} & \textbf{Cifra} & \textbf{Valor} \\
\midrule
T1 & Número de observaciones $n$                             & \marcador{~} \\
T1 & Suma de la serie (dígito de verificación)               & \marcador{~} \\
T1 & Autocorrelación en el rezago 1, $r_1$                   & \marcador{~} \\
T1 & Autocorrelación en el rezago 6, $r_6$                   & \marcador{~} \\
T1 & ADF: estadístico                                        & \marcador{~} \\
T1 & ADF: valor $p$                                          & \marcador{~} \\
T1 & KPSS con nivel: estadístico y valor $p$                 & \marcador{~} \\
T1 & KPSS con tendencia: estadístico y valor $p$             & \marcador{~} \\
T1 & \texttt{ndiffs()}                                       & \marcador{~} \\
T1 & Varianza de $y_t$ y de $\nabla y_t$                     & \marcador{~} \\
\midrule
T2 & $F_T$ y $F_S$ con la descomposición \textbf{recibida}   & \marcador{~} \\
T2 & $F_T$ y $F_S$ ya \textbf{corregida}                     & \marcador{~} \\
T2 & Índice estacional de julio, en el modelo recibido       & \marcador{~} \\
T2 & ACF del residual en el rezago 12, $r_{12}$              & \marcador{~} \\
T2 & Banda del correlograma, $\pm 1.96/\sqrt{n}$             & \marcador{~} \\
T2 & ¿Cruza $r_{12}$ la banda? Antes / después de corregir   & \marcador{~} \\
\midrule
T3 & El orden $(d, D)$ correcto, y con qué transformación     & \marcador{~} \\
T3 & Varianza mínima \emph{dentro} de esa transformación     & \marcador{~} \\
\midrule
T4 & \texttt{nsdiffs()} con la frecuencia mal declarada      & \marcador{~} \\
T4 & \texttt{nsdiffs()} con la frecuencia correcta           & \marcador{~} \\
T4 & Rezago del pico de la ACF: mal declarada / correcta     & \marcador{~} \\
\bottomrule
\end{tabularx}
\endgroup

\vspace{0.45cm}
\noindent\textbf{T3 · las dos tablas.} Copia aquí la tabla que declares \textbf{correcta}, y señala
en la otra la fila donde su aritmética se delata.

\vspace{0.25cm}
\begingroup\setlength{\parskip}{0pt}
\begin{tabularx}{\textwidth}{@{}l c c r X@{}}
\toprule
\textbf{Transformación} & \textbf{$d$} & \textbf{$D$} & \textbf{Varianza} & \textbf{¿Comparable con la fila de arriba?} \\
\midrule
\marcador{~} & \marcador{~} & \marcador{~} & \marcador{~} & \marcador{~} \\
\marcador{~} & \marcador{~} & \marcador{~} & \marcador{~} & \marcador{~} \\
\marcador{~} & \marcador{~} & \marcador{~} & \marcador{~} & \marcador{~} \\
\bottomrule
\end{tabularx}
\endgroup

\vspace{0.2cm}
{\small\color{grismarcador} T5 no lleva cifras propias: sus datos vienen dados en el enunciado y son
los mismos para todo el curso.}

\newpage

% ============================================================================
% LAS CINCO TAREAS
% ============================================================================
\tarea{T1}{La serie que no se deja clasificar}{20}{700}
\literal{a}{Si tu serie es no estacionaria y de qué manera, justificado con las \textbf{dos} pruebas y con el gráfico. \texttt{ndiffs()} no es una justificación: es una cifra que hay que explicar.}
\marcador{tu respuesta}

\literal{b}{Qué \emph{no} prueba tu resultado. Contesta con el poder de la prueba y con tu $n$, no con el valor $p$.}
\marcador{tu respuesta}

\literal{c}{Un compañero tiene un ADF casi idéntico al tuyo sobre una serie que no se parece a la tuya. ¿Significan lo mismo los dos estadísticos?}
\marcador{tu respuesta}

\literal{d}{Si tus tres evidencias —gráfico, ADF y KPSS— no apuntan al mismo sitio, dilo y decide igual. Una contradicción explicada vale más que un acuerdo fabricado.}
\marcador{tu respuesta}

\tarea{T2}{La descomposición que el gráfico desmiente}{20}{700}
\literal{a}{Qué síntoma delata que el modelo de descomposición que recibiste no es el correcto. Nómbralo \textbf{antes} de arreglarlo, y sin usar la palabra «error»: di qué se ve y dónde.}
\marcador{tu respuesta}

\literal{b}{El índice estacional de julio es una constante. Qué afirma esa constante sobre tu serie, y por qué esa afirmación es falsa aquí.}
\marcador{tu respuesta}

\literal{c}{$F_T$ y $F_S$ antes y después de corregir. Cuál de las dos fuerzas se mueve más, y por qué esa y no la otra.}
\marcador{remite a la hoja de cifras, y comenta aquí lo que las cifras no dicen}

\literal{d}{Qué habría pasado aguas abajo si nadie lo nota. Usa la serie desestacionalizada como ejemplo concreto.}
\marcador{tu respuesta}

\tarea{T3}{La auditoría}{25}{800}
\literal{a}{Cuál de las dos tablas está mal, con evidencia \textbf{interna a la tabla}: comprobable con los números que ves y una calculadora, sin recalcular nada sobre la serie.}
\marcador{tu respuesta, y la aritmética que la sostiene}

\literal{b}{Qué elige de verdad esa tabla, si no es el orden de diferenciación.}
\marcador{tu respuesta}

\literal{c}{Cuál es la pareja $(d, D)$ correcta y con qué evidencia de la propia tabla, y en qué escala la lees. La que declaraste incorrecta también contiene esa respuesta: señala dónde.}
\marcador{tu respuesta}

\literal{d}{Por qué este error no se vería mirando el gráfico de la serie.}
\marcador{tu respuesta}

\literal{e}{Si el criterio «elige el $d$ que minimiza la varianza» es entonces inútil, o si el problema es otro. Decide.}
\marcador{tu respuesta}

\tarea{T4}{Declarar no es transformar}{15}{600}
\literal{a}{Qué síntoma lo delata. \textbf{El gráfico estacional y el de rezagos antes que cualquier cifra}: di qué ves y por qué es imposible.}
\marcador{tu respuesta}

\literal{b}{Arréglalo: el componente estacional y \texttt{nsdiffs()} antes y después, con qué función y por qué esa.}
\marcador{tu respuesta}

\literal{c}{Qué habría pasado aguas abajo si nadie lo nota.}
\marcador{tu respuesta}

\literal{d}{Si declarar la frecuencia a mano es tan peligroso, por qué existe: un caso real en que sí es lo correcto.}
\marcador{tu respuesta}

\tarea{T5}{Auditar a la IA}{20}{650}
\literal{a}{Clasifica las tres afirmaciones como correcta, a medias o falsa, con una cifra o un contraejemplo detrás de cada veredicto.}
\marcador{tu respuesta}

\literal{b}{La falsa, reescrita para que quede correcta sin cambiar de tema.}
\marcador{tu reescritura}

\literal{c}{La de «a medias»: dónde se rompe exactamente, qué parte es cierta, cuál no, y qué cifra lo demuestra.}
\marcador{tu respuesta}

\literal{d}{Una frase: qué tipo de pregunta contesta bien un modelo sobre este material y cuál contesta mal, apoyada en lo que te pasó durante el taller.}
\marcador{una frase}

\newpage
% ============================================================================
% DECLARACIÓN DE USO DE IA
% ============================================================================
\section{Declaración de uso de inteligencia artificial}
\markboth{Declaración de uso de IA}{}

El uso de IA en este taller es \textbf{libre y no hay que pedir permiso}. Lo que se califica —el
10\,\% del escrito, dimensión de comunicación y honestidad— es esta declaración: qué consultaste y
\textbf{cómo lo verificaste}. Declarar no penaliza; no declarar, sí.

Escribe una fila por tarea. Si en alguna no consultaste nada, escríbelo: «no consulté» es una
declaración válida y dejarla en blanco no lo es.

\vspace{0.3cm}
\renewcommand{\arraystretch}{1.4}
\begingroup\setlength{\parskip}{0pt}
\begin{tabularx}{\textwidth}{@{}l X X@{}}
\toprule
\textbf{Tarea} & \textbf{Qué consulté y a qué modelo} & \textbf{Cómo verifiqué la respuesta} \\
\midrule
T1 & \marcador{~} & \marcador{~} \\
T2 & \marcador{~} & \marcador{~} \\
T3 & \marcador{~} & \marcador{~} \\
T4 & \marcador{~} & \marcador{~} \\
T5 & \marcador{~} & \marcador{~} \\
\bottomrule
\end{tabularx}
\endgroup

% ============================================================================
% DECISIONES PARA LA DEFENSA
% ============================================================================
\section{Decisiones que sostendré en la defensa}
\markboth{Decisiones para la defensa}{}

La defensa vale el 40\,\%, y hay una regla que le da dientes: \textbf{una decisión entregada por
escrito que no se pueda sostener en la defensa se recalifica}. Escribe entre tres y cinco
decisiones tuyas —no resultados: decisiones— y de qué tarea salen.

\begin{enumerate}[leftmargin=*,itemsep=4pt]
  \item \marcador{decisión} \hfill \textcolor{grismarcador}{\small (tarea \marcador{Tn})}
  \item \marcador{decisión} \hfill \textcolor{grismarcador}{\small (tarea \marcador{Tn})}
  \item \marcador{decisión} \hfill \textcolor{grismarcador}{\small (tarea \marcador{Tn})}
  \item \marcador{opcional} \hfill \textcolor{grismarcador}{\small (tarea \marcador{Tn})}
  \item \marcador{opcional} \hfill \textcolor{grismarcador}{\small (tarea \marcador{Tn})}
\end{enumerate}

\newpage
% ============================================================================
% ANEXO — no cuenta para el limite de paginas
% ============================================================================
\section*{Anexo. Código ejecutado y entorno}
\markboth{Anexo}{}
\addcontentsline{toc}{section}{Anexo}

Pega aquí el código que ejecutaste de verdad para T1, T2 y T4, y la salida de
\texttt{sessionInfo()}. No se califica su elegancia: sirve para saber sobre qué serie trabajaste si
alguna cifra no cuadra.

\begin{lstlisting}[language=R]
# Pega aqui tu codigo
\end{lstlisting}

\begin{lstlisting}
# Pega aqui la salida de sessionInfo()
\end{lstlisting}

\end{document}
